\section{Background and Related Work}
\label{sec:background}

\textbf{Admission control.} Servers protect themselves by queueing or
shedding work: SEDA~\cite{welsh2001seda}, front-end admission control for
e-commerce sites~\cite{elnikety2004gatekeeper}, flash-crowd
detection~\cite{chen2005news}, and controllers that adjust admission from
measured delay, such as DAGOR~\cite{zhou2018dagor} and
Breakwater~\cite{cho2020breakwater}. These act on requests. A waiting room
acts on people: it admits a student, who then sends several requests, and the
student can refresh, give up or come back.

\textbf{Waiting rooms.} WebQ~\cite{doshi2015webq} held users of an overloaded
server in a virtual queue, told them how long to wait and shaped load to the
server's capacity. Commercial rooms offer first-come-first-served and lottery
orders and session lifetimes~\cite{cloudflare2026wr,aws2022vwr}. We are not
aware of a peer-reviewed evaluation of how these settings interact with
students who refresh. \todo{Re-run the venue search in the gap analysis,
section 7, before claiming this.}

\textbf{Retries and metastability.} Retries can keep a system overloaded after
the trigger has gone~\cite{bronson2021metastable,huang2022metastable}; recent
work models how retry policies and capacity loss feed each
other~\cite{metastable2026characterizing,mehan2026retry}. These studies model
automated clients; students are slower, less regular and give up.

\textbf{Wait information.} In a randomised field experiment on a ride-sharing
platform, longer initial wait estimates increased abandonment of a virtual
queue~\cite{yu2023delayinfo}. We test the opposite use of the same effect:
letting students who will not be served in time leave early.

\textbf{Our earlier results.} The first paper~\cite{pathak2026surviving}
introduced ResultBench and the naive and token waiting rooms that serve as
baselines here.
